\documentclass{article}

\usepackage[margin=1in]{geometry}
\usepackage{hyperref}
\usepackage{listings}
\usepackage{graphicx}
\usepackage{booktabs}
\usepackage{enumitem}
\usepackage{xcolor}

\hypersetup{
  colorlinks=true,
  linkcolor=blue,
  urlcolor=blue,
  pdftitle={Swaroop Formulation Industries -- Progressive Website Development Roadmap}
}

\lstset{
  basicstyle=\ttfamily\small,
  breaklines=true,
  frame=single,
  backgroundcolor=\color{gray!10}
}

\title{Swaroop Formulation Industries -- Progressive Website Development Roadmap}
\author{Swaroop Formulation Industries}
\date{2026-03-25}

\begin{document}
\maketitle

\section{Business Context}

\textbf{Swaroop Formulation Industries Pvt. Ltd.} manufactures biodegradable PLA-based plastic bags and bio-medical compostable waste bags in Unnao, Uttar Pradesh, India. The company holds \textbf{ISO 9001:2015} and \textbf{ISO 13485:2016} certifications. It targets \textbf{B2B} logistics, retail, and export markets. Financial projections indicate \textbf{20--45\% ROI} over five years.

\section{Six-Level Roadmap Summary}

\begin{table}[h]
\centering
\begin{tabular}{@{}llp{5.5cm}l@{}}
\toprule
\textbf{Level} & \textbf{Name} & \textbf{Scope} & \textbf{Timeline} \\
\midrule
1 & Beginner & Static site + chatbot; GitHub Pages / Vercel & 1 week \\
2 & Elementary & RAG chatbot with document retrieval; ChromaDB; Render free & 2 weeks \\
3 & Intermediate & Agentic RAG + certifications + basic scraping; LangGraph & 3--4 weeks \\
4 & Upper Intermediate & Adaptive RAG + dashboards + Airflow ETL + PostgreSQL & 4--6 weeks \\
5 & Advanced & Full platform + GenAI automation + CI/CD + Docker & 6--8 weeks \\
6 & Expert & Enterprise K8s + RBAC + observability + IaC & 8--12 weeks \\
\bottomrule
\end{tabular}
\end{table}

\subsection{Level Details}

\begin{itemize}[leftmargin=*]
  \item \textbf{Level 1 (Beginner):} Static site with an embedded chatbot; deploy on GitHub Pages or Vercel; about one week.
  \item \textbf{Level 2 (Elementary):} Retrieval-augmented chatbot using ChromaDB; free tier on Render; about two weeks.
  \item \textbf{Level 3 (Intermediate):} Agentic RAG, certifications showcase, basic web scraping with LangGraph; three to four weeks.
  \item \textbf{Level 4 (Upper Intermediate):} Adaptive RAG, operational dashboards, Airflow ETL, PostgreSQL; four to six weeks.
  \item \textbf{Level 5 (Advanced):} End-to-end platform with GenAI automation, CI/CD, and Docker; six to eight weeks.
  \item \textbf{Level 6 (Expert):} Kubernetes, role-based access control, observability stack, infrastructure as code; eight to twelve weeks.
\end{itemize}

\section{Top-Down Architecture}

\subsection{Frontend (Next.js / React)}

Primary surfaces:

\begin{description}[style=nextline]
  \item[Landing] Marketing and company narrative.
  \item[Chat] Conversational assistant.
  \item[Certs] Certifications and compliance.
  \item[Dash] Analytics and operational views.
  \item[Admin] Configuration and content management.
\end{description}

\subsection{Backend (FastAPI / Python)}

\begin{itemize}[leftmargin=*]
  \item \textbf{REST + WebSocket API} -- synchronous and streaming interactions.
  \item \textbf{Adaptive RAG Engine} -- context-aware retrieval and generation.
  \item \textbf{Agentic Orchestrator} -- multi-step tool use and planning.
  \item \textbf{Airflow ETL} -- scheduled data pipelines.
  \item \textbf{Web Scrapers} -- ingestion from public sources.
\end{itemize}

\subsection{Data Layer}

\begin{tabular}{@{}ll@{}}
\toprule
\textbf{Component} & \textbf{Role} \\
\midrule
ChromaDB / Qdrant & Vector store for embeddings and semantic search \\
PostgreSQL & Relational data, users, metadata \\
S3 / MinIO & Object storage for documents and artifacts \\
Redis & Caching, sessions, rate limiting \\
\bottomrule
\end{tabular}

\section{Version Control Strategy (Progressive)}

Development maturity maps to branching models:

\begin{enumerate}[leftmargin=*]
  \item \textbf{Levels 1--2:} Simple \textbf{main} + \textbf{dev}; integrate on \texttt{dev}, release from \texttt{main}.
  \item \textbf{Levels 3--4:} \textbf{GitFlow} -- \texttt{main}, \texttt{develop}, \texttt{feature/*}, \texttt{release/*}, \texttt{hotfix/*}.
  \item \textbf{Levels 5--6:} \textbf{Trunk-based} development with short-lived branches, feature flags, and protected \texttt{main}.
\end{enumerate}

Tags follow \texttt{v\{MAJOR\}.\{MINOR\}.\{PATCH\}-\{level-descriptor\}} (for example \texttt{v0.1.0-l1-chatbot-basic}).

\section{Documentation Convention}

Project documentation is maintained in \textbf{dual format}:

\begin{itemize}[leftmargin=*]
  \item \textbf{Markdown (\texttt{.md})} -- readable in repositories, wikis, and editors.
  \item \textbf{LaTeX (\texttt{.tex})} -- PDF-ready reports, print, and formal deliverables.
\end{itemize}

Content for each topic should stay aligned across both formats; only presentation (headings, code blocks, tables) is adapted to the medium.

\end{document}
